\subsection{Pré-traitement}\label{sec:preprocessing}

Trois problèmes sont traités dans l'ordre : valeurs manquantes, valeurs non numériques,
écarts d'échelle. Chaque étape est un script de \code{preprocessing/}, expliqué dans le
notebook ; le fichier d'origine n'est jamais modifié.

\subsubsection{Gestion des valeurs manquantes}\label{sec:prep-manquantes}

\paragraph{Colonnes écartées.} Les sorties absentes d'au moins trois quarts des lignes
(dureté, FATT, microstructure) et les entrées vides au moins une fois sur deux et
secondaires (Cu, Co, W, B, Sn, As, Sb, surtout à l'état de traces). Ni, Cr, Mo et Nb,
également vides à plus de 50\,\%, sont gardés : ce sont les principaux éléments
d'alliage (figure~\ref{fig:manquants}).

\begin{figure}[!htbp]
  \centering
  \includegraphics[width=0.82\linewidth]{manquants}
  \caption{Part de valeurs manquantes pour chacune des 30 entrées. En gris, les entrées
    écartées.}
  \label{fig:manquants}
\end{figure}

\paragraph{Valeurs déduites.}
\begin{itemize}
  \item \emph{Type de courant} (manquant sur 215 lignes) : une polarité \code{+} ou
    \code{-} implique le courant continu, \code{0} le courant alternatif ; 59 lignes
    complétées.
  \item \emph{Traitement thermique} (manquant sur 13 lignes notées « traitées » d'un
    même article) : nous supposons le traitement des autres soudures de l'article,
    \SI{600}{\celsius} pendant \SI{0,75}{\hour}.
\end{itemize}

\paragraph{Indicateurs de traitement.} Lue comme une quantité, la température ferait de
\SI{250}{\celsius} un demi-revenu. Or le traitement à \SI{250}{\celsius} /
\SI{14}{\hour} (333 lignes) est probablement un dégazage qui modifie peu le métal, alors
qu'un revenu entre 550 et \SI{750}{\celsius} le transforme. Deux entrées 0/1, « soudure
brute » (610 lignes) et « dégazage », permettent au modèle de traiter ces cas à part.

\paragraph{Valeurs complétées.} Les autres valeurs manquantes des entrées sont remplacées
par la médiane, avec pour chaque colonne un \emph{indicateur} 0/1 « valeur manquante » :
l'absence renseigne sur l'article (section~\ref{sec:exploration}). Cet indicateur décrit
la collecte des données, pas le métal (« nickel manquant » ne veut pas dire « sans
nickel ») : il ne servira pas aux recommandations. Les résultats d'essais ne sont jamais
complétés.

\outrouver{colonnes écartées : liste \code{DROPPED} dans \code{build_model_table.py}
(étape 5) ; type de courant : \code{categorical_clean.py} (étape 4) ; traitement inconnu :
\code{heat_treatment.py} (étape 4b) ; indicateurs de traitement : \code{INDICATORS} dans
\code{ml_dataset.py} ; imputation : \code{make_preprocessor} (étape 6).}

\afaire{\item Vérifier dans l'article Gar\&K-1975 le traitement des soudures de
    \SI{25}{\milli\metre} (supposé identique à celui des \SI{12}{\milli\metre}).
  \item Comparer l'imputation par la médiane à une imputation qui tient compte des autres
    colonnes (plus proches voisins, imputation itérative).}

\subsubsection{Gestion des variables non numériques}

\paragraph{Nombres écrits en texte.} Voir le tableau~\ref{tab:decodage}.

\begin{table}[!htbp]
  \centering\small
  \begin{tabularx}{\linewidth}{@{}l l r Y@{}}
    \toprule
    \textbf{Codage} & \textbf{Col.} & \textbf{Valeurs} & \textbf{Règle et justification} \\
    \midrule
    \code{<x} (limite de détection) & 14 col. & 1576 & $x$. La vraie valeur est entre 0
                                                       et $x$. Autre choix possible : $x/2$. \\
    V \code{<5}                     & 9       & 9    & 0,0005\,\%. C'est 5\,ppm : ailleurs,
                                                       V ne dépasse jamais 0,32\,\%. \\
    Ti, Al \code{<0.01}             & 14, 16  & 32   & \SI{100}{ppm}. C'est 0,01\,\% :
                                                       0,01\,ppm serait irréaliste. \\
    \code{67tot33res} (azote)       & 15      & 59   & 67. On garde l'azote total, comme
                                                       sur les autres lignes. \\
    \code{150-200} (entre passes)   & 27      & 38   & \SI{175}{\celsius}, le milieu de
                                                       l'intervalle. \\
    \bottomrule
  \end{tabularx}
  \caption{Conversion des nombres écrits sous forme de texte.}
  \label{tab:decodage}
\end{table}

\paragraph{Variables catégorielles.} Type de courant, polarité et procédé sont convertis
par \emph{encodage one-hot} (une colonne 0/1 par catégorie) : étant nominales, un
encodage ordinal (0, 1, 2\ldots) leur imposerait un ordre qui n'existe pas. Les 10 codes de procédé sont
d'abord regroupés en 5 familles selon la protection du métal fondu (enrobage, flux en
poudre, fil fourré ou gaz), principale différence entre procédés à l'arc
(tableau~\ref{tab:procedes}). Cela fusionne les synonymes (ShMA et MMA) et les codes de 4
à 7 lignes, trop rares pour être appris ; la famille sous gaz reste hétérogène.

\begin{table}[!htbp]
  \centering\small
  \begin{tabular}{@{}l l r@{}}
    \toprule
    \textbf{Famille} & \textbf{Codes d'origine (lignes)} & \textbf{Lignes} \\
    \midrule
    MMA (arc manuel, électrode enrobée) & MMA (1140), ShMA (40, même procédé) & 1180 \\
    SA (arc submergé sous flux)         & SA (261), NGSAW (18), SAA (4)       & 283 \\
    TSA (arc submergé, deux fils)       & TSA (87)                            & 87 \\
    FCA (fil fourré)                    & FCA (87)                            & 87 \\
    GSA (arc sous protection gazeuse)   & NGGMA (7), GMAA (4), GTAA (4)       & 15 \\
    \bottomrule
  \end{tabular}
  \caption{Regroupement des procédés de soudage en cinq familles.}
  \label{tab:procedes}
\end{table}

\outrouver{\code{decensor.py} (étape 1), \code{nitrogen_total.py} (étape 2),
\code{interpass_range.py} (étape 3), \code{categorical_clean.py} (étape 4) ; encodage :
\code{make_preprocessor} dans \code{ml_dataset.py} (étape 6).}

\afaire{\item Tests de sensibilité : \code{<x} $\rightarrow x$ ou $x/2$ ; azote total ou
    total + résiduel.}

\subsubsection{Mise à l'échelle}

Le carbone varie entre 0,03 et 0,18\,\%, le courant entre 115 et \SI{900}{\ampere}. Nous
\emph{standardisons} toutes les variables numériques (moyenne 0, écart-type 1) :
\begin{itemize}
  \item \textbf{nécessaire} pour les modèles sensibles à l'échelle : Ridge, Lasso, ACP,
    PLS, plus proches voisins, SVR, réseaux de neurones ;
  \item \textbf{sans effet} sur les prédictions des autres : régression linéaire simple,
    arbres et forêts (un arbre coupe au même endroit quelle que soit l'unité).
\end{itemize}
Tous les modèles reçoivent ainsi la même préparation. Limites : la standardisation est
sensible aux valeurs extrêmes et ne corrige pas l'asymétrie ; une mise à l'échelle
robuste ou une transformation logarithmique seront comparées.

\paragraph{Éviter la fuite d'information.} Imputation, encodage et standardisation
utilisent des statistiques des données (médiane, moyenne\ldots). Calculées sur toute la
base, elles feraient profiter l'entraînement des données de test et rendraient les
résultats trop optimistes. Un \code{Pipeline} scikit-learn les recalcule sur les seules
données d'entraînement, à chaque étape de l'évaluation.

\outrouver{\code{make_preprocessor} dans \code{preprocessing/ml_dataset.py} (étape 6).}

\afaire{\item Transformations des variables asymétriques (log1p, mise à l'échelle robuste).
  \item Vérifier la température Charpy de \SI{188}{\celsius} (\code{Ditt-0.5rch2}).}
